\begin{thebibliography}{37}
\bibitem{r1} James Aisenberg, Maria Luisa Bonet, and Sam Buss. 2-d tucker is PPA complete. Electronic Colloquium on Computational Complexity (ECCC), 22:163, 2015. URL: http://eccc.hpi-web. de/report/2015/163.
\bibitem{r2} Noga Alon. Splitting necklaces. Advances in Mathematics, 63(3):247–253, 1987. doi:10.1016/ 0001-8708(87)90055-7.
\bibitem{r3} Noga Alon, Shmuel Friedland, and Gil Kalai. Regular subgraphs of almost regular graphs. Journal of Combinatorial Theory, Series B, 37(1):79–91, 1984.
\bibitem{r4} Noga Alon and Douglas B. West. The Borsuk-Ulam theorem and bisection of necklaces. Proceedings of the American Mathematical Society, 98(4):623–628, 1986. doi: 10.1090/S0002-9939-1986-0861764-9.
\bibitem{r5} Paul Beame, Stephen A. Cook, Jeff Edmonds, Russell Impagliazzo, and Toniann Pitassi. The relative complexity of NP search problems. J. Comput. Syst. Sci., 57(1):3–19, 1998. doi:10.1006/jcss.1998.1575.
\bibitem{r6} Paul Beame and Søren Riis. More on the relative strength of counting principles. In Proceedings of the DIMACS Workshop on Proof Complexity and Feasible Arithmetics, volume 39, pages 13–35, 1998.
\bibitem{r7} Richard Beigel and John Gill. Counting classes: Thresholds, parity, mods, and fewness. Theor. Comput. Sci., 103(1):3–23, 1992. doi:10.1016/0304-3975(92)90084-S.
\bibitem{r8} Aleksandrs Belovs, Gábor Ivanyos, Youming Qiao, Miklos Santha, and Siyi Yang. On the polynomial parity argument complexity of the combinatorial nullstellensatz. In 32nd Computational Complexity Conference, CCC 2017, July 6-9, 2017, Riga, Latvia, pages 30:1– 30:24, 2017. doi:10.4230/LIPIcs.CCC.2017.30.
\bibitem{r9} Nir Bitansky, Omer Paneth, and Alon Rosen. On the cryptographic hardness of finding a nash equilibrium. In 2015 IEEE 56th Annual Symposium on Foundations of Computer Science, pages 1480–1498. IEEE, 2015.
\bibitem{r10} Josh Buresh-Oppenheim and Tsuyoshi Morioka. Relativized NP search problems and propositional proof systems. In 19th Annual IEEE Conference on Computational Complexity (CCC 2004), 21-24 June 2004, Amherst, MA, USA, pages 54–67, 2004. doi: 10.1109/CCC.2004.1313795.
\bibitem{r11} Joshua Buresh-Oppenheim. On the TFNP complexity of factoring. Manuscript, 2006. URL: http://www.cs.toronto.edu/~bureshop/factor.pdf.
\bibitem{r12} Samuel R. Buss, Dima Grigoriev, Russell Impagliazzo, and Toniann Pitassi. Linear gaps between degrees for the polynomial calculus modulo distinct primes. J. Comput. Syst. Sci., 62(2):267–289, 2001. doi:10.1006/jcss.2000.1726.
\bibitem{r13} Samuel R. Buss and Alan S. Johnson. Propositional proofs and reductions between NP search problems. Ann. Pure Appl. Logic, 163(9):1163–1182, 2012. doi:10.1016/j.apal.2012.01.015.
\bibitem{r14} I. Bárány, S. B. Shlosman, and A. Szücs. On a topological generalization of a theorem of tverberg. Journal of the London Mathematical Society, s2-23(1):158–164, 1981. doi: 10.1112/jlms/s2-23.1.158.
\bibitem{r15} Claude Chevalley. Démonstration d’une hypothèse de m. artin. Abhandlungen aus dem Mathematischen Seminar der Universität Hamburg, 11(1):73–75, December 1935. doi:10. 1007/BF02940714.
\bibitem{r16} Arka Rai Choudhuri, Pavel Hubácek, Chethan Kamath, Krzysztof Pietrzak, Alon Rosen, and Guy N Rothblum. Finding a nash equilibrium is no easier than breaking fiat-shamir. In Proceedings of the 51st Annual ACM SIGACT Symposium on Theory of Computing, pages 1103–1114. ACM, 2019.
\bibitem{r17} Constantinos Daskalakis, Paul W. Goldberg, and Christos H. Papadimitriou. The complexity of computing a nash equilibrium. SIAM J. Comput., 39(1):195–259, 2009. doi:10.1137/ 070699652.
\bibitem{r18} Constantinos Daskalakis and Christos H. Papadimitriou. Continuous local search. In Proceedings of the Twenty-Second Annual ACM-SIAM Symposium on Discrete Algorithms, SODA 2011, San Francisco, California, USA, January 23-25, 2011, pages 790–804, 2011. doi:10.1137/1.9781611973082.62.
\bibitem{r19} Xiaotie Deng, Jack R. Edmonds, Zhe Feng, Zhengyang Liu, Qi Qi, and Zeying Xu. Understanding PPA-Completeness. In Ran Raz, editor, 31st Conference on Computational Complexity (CCC 2016), volume 50 of Leibniz International Proceedings in Informatics (LIPIcs), pages 23:1–23:25, Dagstuhl, Germany, 2016. Schloss Dagstuhl–Leibniz-Zentrum fuer Informatik. doi:10.4230/LIPIcs.CCC.2016.23.
\bibitem{r20} Aris Filos-Ratsikas and Paul W. Goldberg. Consensus halving is ppa-complete. In Proceedings of the 50th Annual ACM SIGACT Symposium on Theory of Computing, STOC 2018, Los Angeles, CA, USA, June 25-29, 2018, pages 51–64, 2018. doi:10.1145/3188745.3188880.
\bibitem{r21} Aris Filos-Ratsikas and Paul W. Goldberg. The complexity of splitting necklaces and bisecting ham sandwiches. In STOC (to appear), 2019. arXiv:1805.12559.
\bibitem{r22} C. Goldberg and D. West. Bisection of circle colorings. SIAM Journal on Algebraic Discrete Methods, 6(1):93–106, 1985. doi:10.1137/0606010.
\bibitem{r23} Mika Göös, Pritish Kamath, Robert Robere, and Dmitry Sokolov. Adventures in monotone complexity and TFNP. In 10th Innovations in Theoretical Computer Science Conference, ITCS 2019, January 10-12, 2019, San Diego, California, USA, pages 38:1–38:19, 2019. doi:10.4230/LIPIcs.ITCS.2019.38.
\bibitem{r24} Michelangelo Grigni. A sperner lemma complete for ppa. Information Processing Letters, 77(5-6):255–259, 2001.
\bibitem{r25} Alexandros Hollender. The classes PPA-k: Existence from arguments modulo k. In Ioannis Caragiannis, Vahab Mirrokni, and Evdokia Nikolova, editors, Web and Internet Economics, pages 214–227, Cham, 2019. Springer International Publishing.
\bibitem{r26} Emil Jerábek. Integer factoring and modular square roots. J. Comput. Syst. Sci., 82(2):380–394, 2016. doi:10.1016/j.jcss.2015.08.001.
\bibitem{r27} Alan S. Johnson. Reductions and propositional proofs for total NP search problems. UC San Diego Electronic Theses and Dissertations, 2011. URL: https://escholarship.org/uc/ item/89r774x7.
\bibitem{r28} David S. Johnson, Christos H. Papadimitriou, and Mihalis Yannakakis. How easy is local search? J. Comput. Syst. Sci., 37(1):79–100, 1988. doi:10.1016/0022-0000(88)90046-3.
\bibitem{r29} Ilan Komargodski, Moni Naor, and Eylon Yogev. White-box vs. black-box complexity of search problems: Ramsey and graph property testing. Journal of the ACM (JACM), 66(5):34, 2019.
\bibitem{r30} Pravesh K Kothari and Ruta Mehta. Sum-of-squares meets nash: lower bounds for finding any equilibrium. In Proceedings of the 50th Annual ACM SIGACT Symposium on Theory of Computing, pages 1241–1248. ACM, 2018.
\bibitem{r31} Donald L. Kreher and Douglas R. Stinson. Combinatorial Algorithms: Generation, Enumeration, and Search, volume 7 of Discrete Mathematics and Its Applications. CRC Press, 1998.
\bibitem{r32} Nimrod Megiddo and Christos H. Papadimitriou. On total functions, existence theorems and computational complexity. Theor. Comput. Sci., 81(2):317–324, 1991. doi:10.1016/ 0304-3975(91)90200-L.
\bibitem{r33} Christos H. Papadimitriou. On the complexity of the parity argument and other inefficient proofs of existence. J. Comput. Syst. Sci., 48(3):498–532, 1994. doi:10.1016/S0022-0000(05) 80063-7.
\bibitem{r34} Christian Reiher. On kemnitz’conjecture concerning lattice-points in the plane. The Ramanujan Journal, 13(1-3):333–337, 2007.
\bibitem{r35} Aviad Rubinstein. Settling the complexity of computing approximate two-player nash equilibria. In IEEE 57th Annual Symposium on Foundations of Computer Science, FOCS 2016, 9-11 October 2016, Hyatt Regency, New Brunswick, New Jersey, USA, pages 258–265, 2016.
\bibitem{r36} Katerina Sotiraki, Manolis Zampetakis, and Giorgos Zirdelis. Ppp-completeness with connections to cryptography. In 59th IEEE Annual Symposium on Foundations of Computer Science, FOCS 2018, Paris, France, October 7-9, 2018, pages 148–158, 2018. doi: 10.1109/FOCS.2018.00023.
\bibitem{r37} Ewald Warning. Bemerkung zur vorstehenden arbeit von herrn chevalley. Abh. Math. Sem. Univ. Hamburg, 11:76–83, 1936.
\end{thebibliography}
